
Athugum að strikin $O_1A$, $AO_2$, $BO_1$ og $BO_2$ eru öll geislar, svo
allar hliðar $O_1A,AO_2,BO_1,BO_2$ eru jafnar $r$. Ferhyrningurinn
$O_1AO_2B$ er því samsíðungur.

Látum $\theta=\angle AO_1B = \angle AO_2B$). Þá má skrifa
flatarmálið sem summa tveggja eins þríhyrninga
\[
S=[O_1AB]+[O_2AB]
  = \tfrac12 r^2\sin\theta+\tfrac12 r^2\sin\theta
  = r^2\sin\theta.
\]
Þetta hámarkast þegar $\theta=90^\circ$, því $\sin\theta\le 1$.
Þá er $S_{\max}=r^2$ og samsíðungurinn verður rétthyrndur; þar sem allar
hliðar eru $r$ er hann í raun ferníngur.

\textit{Önnur lausn (án hornafalla).}
Látum $O_1O_2=2d$ og $M$ vera miðpunkt strengsins $AB$. Þá er $AB$
hornrétt á $O_1O_2$ og $O_1M=O_2M=d$. Í rétthyrnda þríhyrningnum
$O_1MA$ fæst
\[
AM=\sqrt{r^2-d^2}.
\]
Hæð samsíðungsins $O_1AO_2B$ að grunninum $O_1O_2$ er $AM$, svo
\[
S=O_1O_2\cdot AM
  = 2d\sqrt{r^2-d^2}.
\]
Setjum $x=d^2$. Þá er
\[
S^2=4d^2(r^2-d^2)=4r^2x-4x^2
= -4\!\left(x-\frac{r^2}{2}\right)^2+r^4 \le r^4.
\]
Því er $S\le r^2$ og hámarkið $S_{\max}=r^2$ næst þegar
$x=\tfrac{r^2}{2}$, þ.e. $d=\tfrac{r}{\sqrt2}$ (svo $O_1O_2=\sqrt2\,r$).
Í því tilfelli eru horn samsíðungsins rétthyrnd og allir leggir hans eru
$r$, þannig að hann er ferníngur með flatarmálið $r^2$.